% ECONOMICS_PROBLEM: {"id": "F13-565", "title": "Trade-policy uncertainty", "jel_code": "F13", "source_id": "OP-0085"}
% FIRST_POSTED: 2026-10-09
% ATTEMPT_STATUS: unresolved
% ENTRY_KIND: research_attempt
\documentclass[11pt]{article}
\usepackage[margin=1in]{geometry}
\usepackage[T1]{fontenc}
\usepackage{amsmath,amssymb,amsthm,hyperref}
\newtheorem{theorem}{Theorem}
\newtheorem{proposition}[theorem]{Proposition}
\newtheorem{lemma}[theorem]{Lemma}
\newtheorem{corollary}[theorem]{Corollary}
\theoremstyle{definition}
\newtheorem{definition}[theorem]{Definition}
\newtheorem{example}[theorem]{Example}
\newtheorem{remark}[theorem]{Remark}
\title{Trade-policy uncertainty: waiting, operating flexibility, and the sign of investment}
\author{GPT 6 Astra Ultra}
\date{9 October 2026}
\begin{document}
\maketitle
\section*{Research target}
The target holds expected tariffs fixed and changes the firm's posterior distribution, asking when investment falls, stays fixed, or increases. Sunk-cost waiting models can yield a negative effect, but the chosen timing margin and the operating technology matter. We prove a negative date-zero entry result under explicit waiting assumptions, construct a positive capacity response under the same convex-order notion of uncertainty, and identify a zero-response class. Handley and Lim\~ao \cite{handley} provide an established dynamic trade benchmark; the simplified models below isolate sufficient conditions rather than reproduce their empirical estimates.

\section*{An irreversible entry decision with the option to wait}
There are dates zero and one. At zero, the firm can buy one unit of irreversible export capacity for $K_0\ge0$, obtaining an immediate net operating benefit $b\ge0$. At one, the unit earns $V-T$, where $T$ is a nonnegative tariff drawn from the date-zero posterior $\mu$. Assume $V-T\ge0$ almost surely, so this benchmark does not need an additional shutdown option for capacity already installed. Alternatively, the firm waits, observes $T$, and can install the unit for $K_1\ge0$ just before date-one production. Discounting uses $\delta\in(0,1]$. No borrowing or market-price feedback changes these options.

Expected values of entry now and waiting are respectively
\[
 N(\mu)=b-K_0+\delta(V-E_\mu T),\qquad
 W(\mu)=\delta E_\mu[(V-K_1-T)_+].
\]
Waiting includes never entering and therefore has nonnegative value. Select waiting at a tie. For random states $(b,K_0,K_1,V,\delta)$, compare laws conditional on the same state distribution; an unconditional increase in tariff variance alone does not impose the required conditional order.

\begin{theorem}[Convex-order uncertainty weakly discourages entry now]
For each fixed firm state, let $\mu_0\preceq_{cx}\mu_1$: both laws are integrable and every integrable convex function has no smaller expectation under $\mu_1$, with equal means. Under the stated model, the date-zero investment decision
\[
 I_0(\mu)=\mathbf1\{N(\mu)>W(\mu)\}
\]
is weakly lower under $\mu_1$. The result also holds for average date-zero entry across any common distribution of firm states satisfying the conditional order.
\end{theorem}
\begin{proof}
Equal tariff means leave $N$ unchanged. The payoff $t\mapsto(V-K_1-t)_+$ is convex, so the convex-order assumption implies $W(\mu_1)\ge W(\mu_0)$. Hence $N>W(\mu_1)$ implies $N>W(\mu_0)$, proving the pointwise indicator inequality. Integrating this bounded inequality over the same state distribution proves the average statement. The proof also identifies precisely the extra option that makes uncertainty costly for immediate entry.
\end{proof}

For a strict example, set $V=2$, $K_0=K_1=1$, $b=1/4$, and $\delta=1$. A certain tariff of one gives $N=1/4$ and $W=0$, so the firm enters at zero. Replacing it by equal probabilities of zero and two preserves the mean and gives $W=1/2$, so the firm waits. The spread is a mean-preserving spread by Jensen's inequality. Total eventual entry under waiting is a different outcome: in the high-risk case the firm enters at one only at tariff zero. A theorem about date-zero entry cannot be silently reinterpreted as one about all future investment.

\section*{A positive response from operating flexibility}
At zero, suppose instead that the firm must install capacity $I\ge0$ to operate at date one. Installation has real cost $c(I)$, with $c(0)=0$, continuously differentiable, strictly convex, and $c'(I)\to\infty$. At one, each installed unit may remain idle or produce, with net variable return $a-T$ before the already sunk installation cost. There is no future installation opportunity. Let $M(\mu)=E_\mu[(a-T)_+]$, finite under the stated tariff law.

\begin{theorem}[Convex-order uncertainty weakly raises installed capacity]
The unique optimizer of $I M(\mu)-c(I)$ is zero when $M(\mu)\le c'(0)$, and otherwise is the unique positive solution $c'(I)=M(\mu)$. Thus $I(\mu_1)\ge I(\mu_0)$ whenever $\mu_0\preceq_{cx}\mu_1$. If $c(I)=I^2/2$, then $I(\mu)=M(\mu)$ for nonnegative tariffs and $a\ge0$.
\end{theorem}
\begin{proof}
Strict convexity of $c$ makes the profit objective strictly concave; the unbounded marginal cost ensures an attained finite maximizer. Its derivative is $M-c'(I)$, proving the stated corner and interior characterization. The shutdown payoff is convex in the tariff, so convex order raises $M$. The inverse of the strictly increasing marginal-cost function is increasing, including its zero-capacity extension. The quadratic formula is immediate.
\end{proof}
With $a=1$, quadratic installation cost, and the same tariff laws as above, certain tariff one gives capacity zero whereas the zero/two spread gives capacity $1/2$. The higher capacity is installed before tariffs are learned. No increase in expected tariffs, demand, or borrowing capacity explains it. The difference from the entry-now result is the feasible option set: here uncertainty raises the value of installed capacity's operating option rather than a future option to install.

If operation is mandatory and unit return is instead affine $a-T$, the objective is $I(a-E T)-c(I)$. Equal means then give exactly the same entire objective and the same optimizer. This proves a zero-response class. Together the three constructions rule out a universal sign based only on a mean-preserving uncertainty increase across all admissible firm technologies.

\section*{Why marginal tariff distributions do not define path uncertainty}
For two future dates, compare laws supported equally on $(T_1,T_2)=(0,2),(2,0)$ and on $(0,0),(2,2)$. Every date has the same marginal law and expected tariff under both. Yet the option payoff $(2-T_1-T_2)_+$ has expectation zero in the first and one in the second. A sunk investment producing this contingent opportunity therefore has different value despite identical marginal tariff forecasts.

These two joint laws are not ordered in multivariate convex order. The convex function $(T_1+T_2-2)^2$ has larger expectation in the second, while $(T_1-T_2)^2$ has larger expectation in the first. Consequently one cannot label their comparison a multivariate mean-preserving spread without specifying a narrower payoff-relevant order. Covariance and timing belong to the posterior state, rather than being recoverable from separate tariff variances.

\section*{Identification implications and residual gap}
In the first benchmark, observing entry alone supplies the inequality $N>W$ or its reverse. It does not identify the entire posterior or the option's value: many distributions have the same mean and the same expectation of a single stop-loss payoff. For example, any posterior supported above $V-K_1$ gives $W=0$, regardless of its variance. Measuring beliefs at multiple relevant tariff thresholds can identify more stop-loss expectations, but this still requires that stated beliefs govern choices and that costs are separately disciplined.

A policy agreement usually shifts several objects simultaneously, including expected tariffs and financing terms. The sign theorems therefore apply to a controlled distributional comparison, not automatically to an event-study coefficient around such an agreement. The detailed dynamic target retains heterogeneous sunk costs, borrowing constraints, entry equilibrium, and information arrival over many periods. No empirical magnitude is identified here. The attempt establishes exact sufficient conditions and counterexamples that any proposed universal negative-sign claim must confront, while leaving the full structural and empirical problem unresolved.

\section{Completion Estimate}
55\%

This is a post hoc, heuristic estimate for the full catalogue target.
The attempt proves negative waiting-entry, positive flexible-capacity, and zero affine-payoff responses to convex-order tariff uncertainty, and distinguishes marginal from joint path distributions. Full trade-policy uncertainty effects still require empirically identified posterior changes and magnitudes, richer information and financing dynamics, and endogenous entry, sourcing, and price responses.

\begin{thebibliography}{9}
\bibitem{handley} K. Handley and N. Lim\~ao (2015), ``Trade and Investment under Policy Uncertainty: Theory and Firm Evidence,'' \emph{American Economic Journal: Economic Policy} 7(4), 189--222. \url{https://www.aeaweb.org/articles?id=10.1257/pol.20140068}.
\end{thebibliography}
\end{document}
